\clearpage
\VerdierPhysicalPageAnchor{0152}%
\begingroup
\small
\setlength{\parskip}{0pt}%
\setlength{\emergencystretch}{1.5em}%

\VerdierSemanticAnchor{ch3}
\section{Abélianisation des catégories triangulées.}

\VerdierSemanticAnchor{ch3-3-1}
\subsection{Construction de la catégorie
\texorpdfstring{$\mathrm A(\mathcal D)$}{A(D)}.}

\VerdierSemanticAnchor{ch3-3-1-1}
\noindent\textbf{3.1.1.}\quad
Soient $\mathcal D$ une catégorie additive et
$\mathcal{F}\ell_1(\mathcal D)$ la catégorie des morphismes de $\mathcal D$
(chap. I, 3.3.6). Un morphisme

\VerdierSemanticAnchor{ch3-3-1-1-square}
\[
\begin{tikzcd}[column sep=4.2em,row sep=2.5em]
X \arrow[r,"f"] \arrow[d,"g"'] & Y \arrow[d,"h"] \\
X' \arrow[r,"f'"'] & Y'
\end{tikzcd}
\]

\VerdierSemanticAnchor{ch3-3-1-1-negligible-morphism}
\noindent de $\mathcal{F}\ell_1(\mathcal D)$ est dit \emph{négligeable} si
$f'g=hf=0$. Le composé, dans les deux sens, d'un morphisme négligeable et
d'un morphisme quelconque est un morphisme négligeable ; la somme directe de
deux morphismes négligeables est un morphisme négligeable.

\VerdierSemanticAnchor{ch3-3-1-1-quotient-category}
Désignons par $\mathrm A(\mathcal D)$ la catégorie dont les objets sont ceux de
$\mathcal{F}\ell_1(\mathcal D)$ et dont les morphismes sont ceux de
$\mathcal{F}\ell_1(\mathcal D)$ à morphismes négligeables près ; autrement
dit, pour tous objets $U,V$ de $\mathrm A(\mathcal D)$, le groupe
$\operatorname{Hom}_{\mathrm A(\mathcal D)}(U,V)$ est le quotient de
$\operatorname{Hom}_{\mathcal{F}\ell_1(\mathcal D)}(U,V)$ par le sous-groupe
des morphismes négligeables. La catégorie $\mathrm A(\mathcal D)$ est additive.

\VerdierSemanticAnchor{ch3-3-1-1-canonical-quotient}
Soit $C$ une catégorie ; la composition avec le foncteur canonique
\[
\mathcal{F}\ell_1(\mathcal D)\longrightarrow\mathrm A(\mathcal D)
\]
fournit un foncteur

\VerdierSemanticAnchor{ch3-3-1-1-1}
\[
\mathcal H\!om(\mathrm A(\mathcal D),C)
\longrightarrow
\mathcal H\!om(\mathcal{F}\ell_1(\mathcal D),C).
\tag{3.1.1.1}
\]
Ce foncteur est pleinement fidèle et injectif sur les objets. Son image est la
sous-catégorie pleine définie par les foncteurs de
$\mathcal{F}\ell_1(\mathcal D)$ dans $C$ qui transforment en morphismes égaux
les morphismes qui diffèrent par un morphisme négligeable.

\VerdierSemanticAnchor{ch3-3-1-2}
\medskip
\noindent\textbf{3.1.2.}\quad
Le foncteur canonique de $\mathcal D$ dans $\mathcal{F}\ell_1(\mathcal D)$,
$X\longmapsto(X\xrightarrow{\operatorname{id}_X}X)$, composé avec le foncteur
canonique $\mathcal{F}\ell_1(\mathcal D)\longrightarrow\mathrm A(\mathcal D)$,
est un foncteur qu'on note

\VerdierSemanticAnchor{ch3-3-1-2-1}
\[
\mathcal H\colon\mathcal D\longrightarrow\mathrm A(\mathcal D).
\tag{3.1.2.1}
\]
Le foncteur $\mathcal H$ est additif. Il est pleinement fidèle et injectif sur
les objets. Il réalise donc $\mathcal D$ comme sous-catégorie pleine de
$\mathrm A(\mathcal D)$.
\endgroup
